\chapter{Língua Portuguesa}

\section{Compreensão, interpretação e construção de sentido}

Compreender um texto significa reconstruir as informações, relações e pressupostos que ele apresenta. Interpretar exige ir além da simples localização de palavras: o leitor precisa identificar a organização argumentativa, a função dos conectores, a referência dos pronomes, o valor dos tempos verbais e os efeitos de escolhas lexicais e sintáticas.

Em provas de alto nível, a distinção central é entre aquilo que o texto \textbf{afirma}, aquilo que ele \textbf{permite inferir} e aquilo que o leitor apenas \textbf{supõe}.

\subsection{Informação explícita, inferência e extrapolação}

Informação explícita aparece diretamente no enunciado. Inferência é uma conclusão necessária ou fortemente sustentada por marcas textuais. Extrapolação ocorre quando se atribui ao texto informação que não pode ser derivada dele.

\begin{exemplo}[inferência legítima]
Considere:

``Embora o número de processos julgados tenha aumentado, o estoque de processos pendentes permaneceu estável.''

É possível inferir que a entrada de novos processos compensou, aproximadamente, o aumento do número de julgamentos ou que algum outro fator manteve o estoque. Não é possível concluir, apenas com essa frase, que o tribunal se tornou mais eficiente ou menos eficiente.
\end{exemplo}

\subsection{Pressuposto e subentendido}

O \termo{pressuposto} é informação acionada por determinada construção linguística e normalmente preservada mesmo quando a frase é negada. O \termo{subentendido} depende mais fortemente do contexto e da situação comunicativa.

Em ``o sistema voltou a apresentar falhas'', o verbo \emph{voltar} pressupõe que já houve falhas antes. Em ``até o setor X cumpriu a meta'', \emph{até} sugere que esse resultado era menos esperado para X.

Palavras como \emph{ainda, já, apenas, mesmo, até, novamente, sequer} alteram o escopo e podem modificar pressupostos relevantes.

\section{Tipologia e gêneros textuais}

\termo{Tipo textual} corresponde a uma organização linguística predominante. \termo{Gênero textual} é uma forma socialmente reconhecida de comunicação, como relatório, notícia, parecer, ofício, artigo ou edital.

\begin{table}[ht]
\centering
\small
\begin{tabularx}{\textwidth}{l X X}
\toprule
Tipo & Objetivo predominante & Marcas comuns \\
\midrule
Narração & relatar acontecimentos & sequência temporal, mudança de estado, ações. \\
Descrição & caracterizar seres, objetos ou ambientes & propriedades, estados, localização, qualificadores. \\
Exposição & apresentar e explicar ideias & definições, classificações, relações conceituais. \\
Argumentação & defender ponto de vista & tese, razões, evidências, contraposição, conclusão. \\
Injunção & orientar ação & instruções, comandos, procedimentos, imperativos. \\
\bottomrule
\end{tabularx}
\end{table}

Um mesmo gênero pode combinar tipos. Um relatório de auditoria pode conter descrição da situação encontrada, exposição de critérios e argumentação para fundamentar uma conclusão.

\section{Coerência e coesão}

\termo{Coerência} diz respeito à compatibilidade global de sentidos do texto. \termo{Coesão} refere-se aos mecanismos linguísticos que conectam partes do texto.

Um texto pode ser gramaticalmente correto e ainda assim incoerente. Também pode ser coerente em termos globais, mas apresentar problemas de coesão que dificultem identificar referentes e relações lógicas.

\subsection{Coesão referencial}

A referenciação ocorre por pronomes, expressões nominais, elipse, repetição controlada e outros mecanismos.

\begin{exemplo}[ambiguidade pronominal]
``O auditor informou ao gestor que seu relatório continha inconsistências.''

O possessivo \emph{seu} pode retomar \emph{auditor} ou \emph{gestor}. Uma reescrita pode eliminar a ambiguidade:

``O auditor informou ao gestor que o relatório do gestor continha inconsistências.''
\end{exemplo}

\subsection{Anáfora e catáfora}

Na \termo{anáfora}, o elemento retoma informação já apresentada. Na \termo{catáfora}, antecipa informação que será apresentada depois.

``O relatório foi concluído. \textbf{Ele} será encaminhado amanhã.'' — \emph{ele} é anafórico.

``Só lhe peço \textbf{isto}: revise o documento.'' — \emph{isto} antecipa o conteúdo posterior.

\subsection{Coesão sequencial}

Conectores estabelecem relações semânticas entre orações, períodos e parágrafos.

\begin{table}[ht]
\centering
\small
\begin{tabularx}{\textwidth}{l X}
\toprule
Relação & Conectores frequentes \\
\midrule
Adição & e, além disso, bem como, ainda. \\
Adversidade & mas, porém, contudo, entretanto, todavia. \\
Causa & porque, visto que, uma vez que, como. \\
Consequência & de modo que, de forma que, tão... que. \\
Conclusão & portanto, logo, assim, por conseguinte. \\
Concessão & embora, ainda que, mesmo que, apesar de. \\
Condição & se, caso, desde que, contanto que. \\
Finalidade & para que, a fim de que. \\
Comparação & como, tal qual, mais... que, menos... que. \\
\bottomrule
\end{tabularx}
\end{table}

Trocar um conector por outro pode manter a correção sintática e alterar o sentido. ``O sistema falhou, \textbf{portanto} o serviço ficou indisponível'' expressa conclusão. ``O sistema falhou, \textbf{embora} o serviço tenha ficado disponível'' expressa concessão.

\section{Semântica lexical}

As palavras adquirem significado em contexto. Sinonímia absoluta é rara; duas palavras podem ser próximas em determinado uso e incompatíveis em outro.

\subsection{Denotação e conotação}

Denotação corresponde ao emprego mais literal e referencial. Conotação envolve valor figurado, afetivo ou associativo.

``A auditoria identificou uma \textbf{lacuna} no controle'' usa \emph{lacuna} metaforicamente para indicar ausência ou insuficiência.

\subsection{Polissemia, homonímia e paronímia}

\termo{Polissemia} ocorre quando uma mesma palavra possui sentidos relacionados. \termo{Homonímia} envolve formas iguais ou próximas com significados distintos. \termo{Parônimos} são palavras semelhantes na forma, mas diferentes no significado.

Pares frequentes:
\begin{itemize}
\item ratificar = confirmar; retificar = corrigir;
\item descrição = ato de descrever; discrição = reserva;
\item iminente = prestes a ocorrer; eminente = elevado, destacado;
\item deferir = conceder; diferir = adiar ou distinguir.
\end{itemize}

\section{Ortografia oficial}

Ortografia abrange grafia das palavras, acentuação, emprego do hífen e formas problemáticas. Em provas, costuma aparecer integrada a reescritas.

\subsection{Acentuação gráfica}

A identificação da sílaba tônica permite classificar palavras em oxítonas, paroxítonas e proparoxítonas. Todas as proparoxítonas são acentuadas. As regras de oxítonas e paroxítonas dependem das terminações previstas pela norma.

Hiatos também geram casos relevantes, como em \emph{saída} e \emph{país}, observadas as exceções normativas.

\subsection{Por que, porque, por quê e porquê}

\begin{table}[ht]
\centering
\small
\begin{tabularx}{\textwidth}{l X}
\toprule
Forma & Uso típico \\
\midrule
por que & pergunta direta/indireta ou sequência preposição + pronome. \\
porque & conjunção explicativa ou causal. \\
por quê & forma tônica no fim de enunciado ou antes de pausa forte. \\
porquê & substantivo, normalmente acompanhado de determinante. \\
\bottomrule
\end{tabularx}
\end{table}

Exemplos: ``Não sei \textbf{por que} o sistema falhou.''; ``Falhou \textbf{porque} estava desatualizado.''; ``O sistema falhou \textbf{por quê}?''; ``O relatório explicou \textbf{o porquê} da falha.''

\subsection{Há e a}

\emph{Há} pode indicar existência ou tempo decorrido: ``há três anos''. \emph{A} pode indicar distância ou tempo futuro: ``daqui a três anos''.

\section{Classes de palavras em funcionamento}

A classificação morfológica precisa ser analisada dentro da frase.

\subsection{Substantivo, adjetivo e advérbio}

Substantivos nomeiam seres, entidades, processos ou conceitos. Adjetivos caracterizam substantivos. Advérbios modificam verbos, adjetivos, outros advérbios ou enunciados inteiros.

Em ``a análise \textbf{técnica} ocorreu \textbf{rapidamente}'', \emph{técnica} caracteriza o substantivo \emph{análise}; \emph{rapidamente} modifica o verbo \emph{ocorreu}.

\subsection{Pronome relativo e conjunção integrante}

Compare:

``O relatório \textbf{que} foi publicado contém ressalvas.''

O \emph{que} retoma \emph{relatório} e exerce função sintática na oração subordinada; é pronome relativo.

``A equipe afirmou \textbf{que} o relatório foi publicado.''

Nesse caso, \emph{que} não retoma antecedente e introduz oração substantiva; é conjunção integrante.

\subsection{Valores de ``se''}

O termo \emph{se} pode exercer funções distintas:
\begin{itemize}
\item pronome reflexivo: ``o servidor se feriu'';
\item pronome apassivador: ``publicaram-se os resultados'';
\item índice de indeterminação do sujeito: ``precisa-se de técnicos'';
\item conjunção condicional: ``se houver recurso, o processo será revisto'';
\item conjunção integrante: ``não se sabe se haverá recurso''.
\end{itemize}

A análise da transitividade verbal é decisiva para diferenciar pronome apassivador de índice de indeterminação.

\section{Estrutura sintática da oração}

A oração organiza-se em torno de um verbo ou locução verbal. O predicado contém o verbo e seus complementos, e o sujeito pode ser expresso, oculto, indeterminado ou inexistente em construções impessoais.

\subsection{Sujeito e núcleo}

O verbo concorda, como regra, com o núcleo do sujeito. Identificar o termo mais próximo do verbo não basta.

``A relação de contratos pendentes \textbf{foi} atualizada.''

O núcleo do sujeito é \emph{relação}, singular; \emph{contratos} integra complemento nominal.

\subsection{Objeto direto e indireto}

Objeto direto liga-se ao verbo sem preposição exigida. Objeto indireto exige preposição determinada pela regência.

``A equipe analisou \textbf{o processo}.'' — objeto direto.

``A equipe obedeceu \textbf{à norma}.'' — objeto indireto introduzido por \emph{a}.

\subsection{Predicativo}

Predicativo atribui característica ao sujeito ou ao objeto:

``O relatório permaneceu \textbf{incompleto}.'' — predicativo do sujeito.

``A comissão considerou o relatório \textbf{incompleto}.'' — predicativo do objeto.

\section{Coordenação e subordinação}

Orações coordenadas possuem independência sintática relativa. Orações subordinadas exercem função sintática em relação a outra oração.

\subsection{Orações subordinadas substantivas}

Podem exercer funções de sujeito, objeto, complemento nominal, predicativo ou aposto.

``É necessário \textbf{que o processo seja revisado}.'' — a oração destacada exerce função de sujeito da estrutura ``é necessário''.

\subsection{Orações subordinadas adjetivas}

Modificam substantivos e são introduzidas, em regra, por pronome relativo.

Compare:

``Os servidores \textbf{que concluíram o curso} receberam certificado.''

A oração é restritiva: seleciona parte do conjunto.

``Os servidores, \textbf{que concluíram o curso}, receberam certificado.''

A oração explicativa apresenta informação acessória e sugere que o conjunto referido já está determinado. A pontuação modifica o alcance semântico.

\subsection{Orações subordinadas adverbiais}

Expressam circunstâncias como causa, condição, concessão, finalidade, consequência e tempo.

``\textbf{Embora houvesse falhas}, o serviço permaneceu disponível.'' — relação concessiva.

``\textbf{Se houver falhas}, o serviço será interrompido.'' — relação condicional.

\section{Pontuação}

Pontuação organiza a estrutura sintática e informacional do texto. A vírgula não corresponde simplesmente a uma pausa de leitura.

\subsection{Casos centrais de vírgula}

Em regra, não se separa por vírgula:
\begin{itemize}
\item sujeito e verbo;
\item verbo e complemento;
\item nome e complemento nominal.
\end{itemize}

A vírgula pode:
\begin{itemize}
\item isolar vocativo e aposto explicativo;
\item marcar expressões intercaladas;
\item separar orações coordenadas;
\item destacar oração adverbial deslocada;
\item delimitar oração adjetiva explicativa;
\item organizar enumerações.
\end{itemize}

\begin{exemplo}[mudança de sentido]
``Os processos que apresentam falhas serão revisados.''

Somente os processos com falhas serão revisados.

``Os processos, que apresentam falhas, serão revisados.''

A construção explicativa atribui a condição ao conjunto já identificado, produzindo leitura diferente.
\end{exemplo}

\subsection{Dois-pontos e ponto e vírgula}

Dois-pontos podem introduzir enumeração, explicação, consequência explicitada ou citação. Ponto e vírgula separa unidades maiores que já contêm vírgulas ou organiza itens complexos.

\section{Concordância verbal}

\subsection{Regra geral}

O verbo concorda com o núcleo do sujeito em número e pessoa.

``Os relatórios \textbf{foram} publicados.''

\subsection{Expressões partitivas}

Com expressões como \emph{a maioria de, grande parte de, metade de}, podem existir possibilidades de concordância conforme a construção e a norma admitida:

``A maioria dos servidores compareceu'' / ``A maioria dos servidores compareceram''.

Em prova, é necessário observar se a questão afirma exclusividade onde a norma admite variação.

\subsection{Haver impessoal}

\emph{Haver} com sentido de existir é impessoal e permanece na terceira pessoa do singular:

``\textbf{Há} inconsistências.''

``\textbf{Deve haver} inconsistências.''

O verbo auxiliar da locução acompanha a impessoalidade.

\subsection{Fazer indicando tempo}

``\textbf{Faz} três anos que o sistema foi implantado.''

Também é impessoal nessa acepção.

\subsection{Existir}

\emph{Existir} possui sujeito e concorda normalmente:

``\textbf{Existem} inconsistências.''

``\textbf{Devem existir} inconsistências.''

\section{Concordância nominal}

Artigos, pronomes, numerais e adjetivos ajustam-se ao substantivo segundo regras de gênero e número.

Em estruturas com \emph{é necessário, é bom, é proibido}, a presença de determinante pode alterar a concordância:

``É necessário cautela.''

``É necessária \textbf{a} cautela.''

\section{Regência verbal e nominal}

Regência descreve a relação de dependência entre um termo e seus complementos.

\begin{table}[ht]
\centering
\small
\begin{tabularx}{\textwidth}{l X}
\toprule
Verbo & Regência em uso frequente de prova \\
\midrule
assistir & no sentido de ver, tradicionalmente rege preposição \emph{a}: assistir ao julgamento. \\
obedecer/desobedecer & regem \emph{a}: obedecer à norma. \\
aspirar & no sentido de desejar, rege \emph{a}; no sentido de sorver, é transitivo direto. \\
visar & no sentido de almejar, tradicionalmente rege \emph{a}; em outros sentidos pode ter regência distinta. \\
preferir & estrutura ``preferir X a Y'', sem reforços como ``mais''. \\
implicar & no sentido de acarretar, é transitivo direto na norma tradicional. \\
\bottomrule
\end{tabularx}
\end{table}

A regência interfere diretamente no emprego de pronomes e da crase.

\section{Crase}

Crase é a fusão de dois sons /a/, normalmente representada pelo acento grave em \emph{à/às}. O caso mais frequente é a fusão da preposição \emph{a} exigida pelo termo regente com o artigo feminino \emph{a/as}.

\subsection{Teste de substituição}

``O auditor referiu-se \textbf{à norma}.''

Substituindo \emph{norma} por masculino:

``O auditor referiu-se \textbf{ao regulamento}.''

O aparecimento de \emph{ao} indica presença simultânea de preposição e artigo, justificando \emph{à} na forma feminina.

\subsection{Casos de não ocorrência}

Não há crase quando falta um dos dois elementos da fusão. Em regra, não ocorre antes de verbo, palavra masculina ou artigo indefinido.

``Começou a revisar.'' — antes de verbo.

``Entregou a documentação a um servidor.'' — não há artigo feminino definido.

\subsection{Demonstrativos}

Pode ocorrer crase em \emph{àquele, àquela, àquilo} quando o termo regente exige preposição \emph{a}:

``Referiu-se \textbf{àquela} decisão.''

\section{Colocação pronominal}

Próclise, mesóclise e ênclise dizem respeito à posição dos pronomes oblíquos átonos em relação ao verbo.

\subsection{Próclise}

Elementos atrativos incluem negação, pronomes relativos, certas conjunções subordinativas e advérbios em determinadas condições:

``Não \textbf{se} identificaram falhas.''

``O relatório que \textbf{se} publicou ontem...''

\subsection{Ênclise}

É frequente quando o verbo inicia a oração e não há fator de próclise:

``Identificaram-\textbf{se} inconsistências.''

\subsection{Mesóclise}

Pode ocorrer com futuro do presente ou futuro do pretérito, em registro formal, quando não há elemento atrativo:

``Far-\textbf{se}-á a revisão.''

\section{Tempos e modos verbais}

O \termo{indicativo} tende a apresentar fatos como assertados ou considerados reais. O \termo{subjuntivo} aparece em contextos de hipótese, possibilidade, desejo, condição, concessão e dependência. O \termo{imperativo} orienta ações.

\subsection{Correlação verbal}

``Se ele \textbf{vier}, analisaremos o processo.''

Futuro do subjuntivo combina com uma consequência projetada.

``Se ele \textbf{viesse}, analisaríamos o processo.''

Imperfeito do subjuntivo combina naturalmente com futuro do pretérito em hipótese contrafactual ou menos provável.

A troca de tempo ou modo pode manter a frase gramatical e alterar certeza, temporalidade ou modalidade.

\section{Vozes verbais}

Na voz ativa, o sujeito é apresentado como agente do processo verbal. Na voz passiva, o sujeito recebe a ação.

``A equipe revisou os dados.''

``Os dados foram revisados pela equipe.''

A transformação preserva o núcleo factual quando realizada corretamente, mas altera a organização informacional: o elemento colocado como sujeito tende a ganhar maior destaque discursivo.

Na passiva sintética:

``Revisaram-se os dados.''

O \emph{se} é pronome apassivador, e o verbo concorda com o sujeito paciente \emph{os dados}.

\section{Reescrita e preservação de sentido}

Questões de reescrita integram sintaxe, semântica e coesão. Uma nova frase pode estar gramaticalmente correta e não ser equivalente à original.

\subsection{Mudança de quantificador}

``Nem todos os contratos foram auditados'' significa que pelo menos um contrato não foi auditado.

``Nenhum contrato foi auditado'' significa que zero contratos foram auditados.

As frases não são equivalentes.

\subsection{Mudança de escopo}

Compare:

``O auditor não afirmou que todos os contratos eram regulares.''

Essa frase nega o ato de afirmar uma proposição universal. Não é equivalente, sem contexto adicional, a:

``O auditor afirmou que todos os contratos eram irregulares.''

Negar uma afirmação não implica afirmar sua contrária.

\subsection{Nominalização}

``A comissão analisou o processo.''

``A análise do processo pela comissão...''

A nominalização converte verbo em nome e pode exigir reajustes de preposição e referência. É comum em textos formais e reescritas.

\begin{cebraspe}[reescrita de alto nível]
Para verificar equivalência, separe quatro perguntas: a estrutura permanece gramatical? os referentes permanecem os mesmos? a relação lógica foi mantida? o grau de certeza, tempo, agente ou quantificador mudou? Uma alternativa pode ser perfeita do ponto de vista gramatical e ainda assim alterar o sentido.
\end{cebraspe}

\section{Síntese conceitual}

\mapafuncional{Língua Portuguesa}
{Sentido textual = léxico + sintaxe + coesão + contexto}
{Interpretação $\rightarrow$ explícito, inferência e pressuposto.\\Coesão $\rightarrow$ referência e conectores.\\Sintaxe $\rightarrow$ termos da oração, subordinação e pontuação.\\Norma-padrão $\rightarrow$ concordância, regência, crase e colocação.\\Reescrita $\rightarrow$ estrutura + sentido.}
{Pronome: identifique o referente.\\Conector: identifique a relação semântica.\\Vírgula: determine a estrutura sintática.\\Verbo: identifique sujeito e regência.\\Reescrita: compare escopo, tempo, agente e quantificador.}
{Frase correta $\neq$ frase equivalente.\\Vírgula $\neq$ pausa livre.\\Sinônimo $\neq$ substituição automática.\\Haver impessoal $\neq$ existir.\\Nem todos $\neq$ nenhum.}
{O que o texto afirma?\\O que apenas permite inferir?\\Quem é o referente?\\Qual é o núcleo do sujeito?\\A preposição é exigida?\\A alteração mudou o alcance da afirmação?}

\begin{resumo}
Língua Portuguesa deve ser estudada como um sistema integrado. Interpretação depende da gramática, e a gramática produz efeitos de sentido. Em provas seletivas, a banca frequentemente altera uma vírgula, um pronome, um conector, uma preposição ou um quantificador e pergunta se a nova redação continua correta e equivalente.
\end{resumo}
